% संपूर्ण मराठी ग्रंथातील प्रकरण 45: स्वयंसिद्धकीय निष्पत्ती
% हा संपादनयोग्य स्रोतखंड आहे. संकलनासाठी संपूर्ण स्रोतसंग्रहातील
% अक्षररूपे, पूर्वभाग, चिन्हव्याख्या आणि आकृती-स्रोत आवश्यक आहेत.
% मूळ: Open Logic Project; मजकूर आणि रूपांतर CC BY 4.0.
% यंत्रानुवाद, दुरुस्ती आणि तपासणी: OpenAI Codex —
% GPT-5.6 Sol आणि GPT-6 Sol, दोन्ही Ultra effort.
% दुरुस्ती-पुनरावलोकन: GPT-6.1 Sol, Ultra effort.
\chapter{स्वयंसिद्धकीय निष्पत्ती}\label{nml:prf:chap}










\section{प्रस्तावना}\label{nml:axs:int:sec}

मूलभूत मोडल भाषेसाठी मोडल प्रतिमानांद्वारे चिन्हार्थमीमांसा
आपल्याकडे आहे. सूत्र वैध असणे, म्हणजे सर्व
प्रतिमानांतील सर्व जगांत सत्य असणे, किंवा प्रतिमानांच्या
किंवा चौकटींच्या एखाद्या वर्गाच्या संदर्भात वैध असणे,
म्हणजे त्या वर्गातील अथवा त्या चौकटीवर आधारित सर्व
प्रतिमानांच्या सर्व जगांत सत्य असणे, या संकल्पनाही
आपल्याकडे आहेत. तर्कशास्त्रात अशा चिन्हार्थविषयक
वैधतेला सहसा निष्पन्नता या सिद्धताविषयक
संकल्पनेशी जोडले जाते. एखाद्या पद्धतीत
निष्पन्नता अशी परिभाषित करायची की
सूत्र वैध असेल तेव्हा आणि केवळ तेव्हाच
निष्पन्न करता येण्याजोगे असेल, हे आपले उद्दिष्ट आहे.

सर्वात सोप्या आणि ऐतिहासिकदृष्ट्या सर्वात जुन्या
निष्पत्ती-पद्धती म्हणजे हिल्बर्ट-प्रकारच्या किंवा
स्वयंसिद्धकीय निष्पत्ती-पद्धती. अनेक मोडल
तर्कशास्त्रांसाठी हिल्बर्ट-प्रकारच्या निष्पत्ती-पद्धती तुलनेने सहज रचता येतात: अधितात्त्विक अभ्यासाच्या
दृष्टीने (उदा., निर्दोषता व संपूर्णता सिद्ध करताना)
त्या सोप्या असतात. पण त्यांमध्ये सूत्रे सिद्ध
करणे, उदाहरणार्थ नैसर्गिक निगमन पद्धतींपेक्षा, बरेच
कठीण असते.

हिल्बर्ट-प्रकारच्या निष्पत्ती-पद्धतीत
सूत्राची निष्पत्ती म्हणजे काही स्वयंसिद्धकांपासून
मोजक्या अनुमान नियमांनी मिळणारी, संबंधित
सूत्र पर्यंत पोचणारी सूत्रांची मालिका.
निष्पत्ती-पद्धत चिन्हार्थमीमांसेशी जुळावी अशी
आपली अपेक्षा असल्याने, निष्पन्न करता येण्याजोगे सूत्रे
सर्व प्रतिमानांत (किंवा स्वयंसिद्धकांचे सर्व आदेशन-दृष्टांत
सर्व जगांत सत्य असलेल्या प्रतिमानांत) सत्य आहेत याची खात्री करावी लागेल.
यासाठी आवश्यक असलेले मोडल तर्कशास्त्रांचे काही
गुणधर्म आधी वेगळे पाहू: ``सामान्य'' मोडल तर्कशास्त्रे.
सामान्य मोडल तर्कशास्त्रांसाठी फक्त दोन अनुमान नियम
गृहीत धरावे लागतात: विध्यात्मक अनुमान आणि
अनिवार्यीकरण. स्वयंसिद्धके म्हणून सर्वतः-सत्य
सूत्रांचे सर्व आदेशन-दृष्टांत आणि संबंधित मोडल
तर्कशास्त्रानुसार काही मोडल स्वयंसिद्धके घेतो.
सर्व प्रतिमानांच्या वर्गातच आपल्याला रस असला,
तरी~$\Ax{K}$ आणि~$\Ax{Dual}$
यांचे सर्व आदेशन-दृष्टांत स्वयंसिद्धके म्हणून घ्यावे
लागतात. एवढ्यानेच न्यूनतम सामान्य मोडल
तर्कशास्त्र~$\Log K$ मिळते.

\begin{defn}
\emph{विध्यात्मक अनुमान} हा नियम पुढील अनुमान-रूपबंध आहे:
\begin{prooftree}
\AxiomC{$\varphi$}
\AxiomC{$\varphi \lif \psi$}
\RightLabel{\MP}
\BinaryInfC{$\psi$}
\end{prooftree}
सूत्र~$\psi$ हे सूत्रे $\varphi$, $\chi$
यांपासून विध्यात्मक अनुमानाने \emph{निष्पन्न होते}
तेव्हा आणि केवळ तेव्हाच, जेव्हा
$\chi \ident \varphi \lif \psi$.
\end{defn}

\begin{defn}
\emph{अनिवार्यीकरण} हा नियम पुढील अनुमान-रूपबंध आहे:
\begin{prooftree}
\AxiomC{$\varphi$}
\RightLabel{\Nec}
\UnaryInfC{$\Box \varphi$}
\end{prooftree}
सूत्र~$\psi$ हे सूत्रे $\varphi$
पासून अनिवार्यीकरणाने निष्पन्न होते तेव्हा आणि
केवळ तेव्हाच, जेव्हा $\psi \ident \Box \varphi$.
\end{defn}

\begin{defn}
स्वयंसिद्धकांच्या संच~$\Sigma$ पासूनची
\emph{निष्पत्ती} म्हणजे सूत्रांची
$\psi_1$, $\psi_2$, \dots, $\psi_n$ अशी मालिका,
ज्यात प्रत्येक $\psi_i$ पुढीलपैकी एक असते:
\begin{enumerate}
\item सर्वतः-सत्य सूत्राचा आदेशन-दृष्टांत; किंवा
\item $\Sigma$ मधील सूत्राचा आदेशन-दृष्टांत; किंवा
\item $j$, $k < i$ असलेल्या दोन सूत्रे
  $\psi_j$, $\psi_k$ यांपासून विध्यात्मक अनुमानाने
  निष्पन्न होणारे; किंवा
\item $j < i$ असलेल्या सूत्र~$\psi_j$
  पासून अनिवार्यीकरणाने निष्पन्न होणारे.
\end{enumerate}
$\psi_n \ident \varphi$ असलेली अशी निष्पत्ती
असल्यास $\varphi$ हे $\Sigma$ पासून
\emph{निष्पन्न करता येण्याजोगे आहे}, असे म्हणतो;
संकेतात $\Sigma \Proves \varphi$.
\end{defn}

वरील परिच्छेदाप्रमाणे $\Sigma$ मध्ये
$\Ax{K}$ आणि~$\Ax{Dual}$ घेतल्यावर,
या व्याख्येनुसार निष्पन्न करता येण्याजोगे सूत्रांचा संच सामान्य मोडल
तर्कशास्त्र बनतो, हे पुढे दिसेल. तसेच, $\Sigma$ मधील
स्वयंसिद्धकांचे सर्व आदेशन-दृष्टांत सर्व जगांत सत्य असलेल्या
प्रत्येक प्रतिमानात प्रत्येक निष्पन्न करता येण्याजोगे सूत्र सत्य असते.
निष्पत्त्यांच्या या गुणधर्माला \emph{निर्दोषता} म्हणतात.
त्याचा व्यत्यास, \emph{संपूर्णता}, सिद्ध करणे अधिक कठीण आहे.












\section{सामान्य मोडल तर्कशास्त्रे}\label{nml:prf:nor:sec}

मोडल सूत्रांचा प्रत्येक संच स्वयंसिद्धकांच्या
एखाद्या संचापासून सिद्ध करता येणाऱ्या
सूत्रांचा संच म्हणून सहज वर्णन करता येत
नाही. मोडल तर्कशास्त्रे सुस्थित असावीत अशी
आपली अपेक्षा आहे. सर्वप्रथम, अभिजात विधानीय
तर्कशास्त्रात आपण निष्पन्न करू शकतो ते सर्व
अजूनही निष्पन्न करता येण्याजोगे असावे; आता सूत्रे
मध्ये $\Box$
    आणि~$\Diamond$ देखील
येऊ शकतात, हे अर्थातच लक्षात घेतले पाहिजे.
यासाठी मोडल तर्कशास्त्रात सर्व सर्वतः-सत्य दृष्टांत
असावेत आणि ते विध्यात्मक अनुमानाखाली बंद असावे,
अशी अट घालतो.

\begin{defn}
  \emph{मोडल तर्कशास्त्र} म्हणजे मोडल सूत्रे
  चा असा संच~$\Sigma$ की, तो
  \begin{enumerate}
  \item सर्व सर्वतः-सत्य सूत्रे सामावतो; आणि
  \item आदेशनाखाली बंद असतो; म्हणजे $\varphi \in \Sigma$
    आणि $\theta_1$, \dots, $\theta_n$ ही सूत्रे
    असतील, तर
    \[    \SSubst{\varphi}{\subst{\theta_1}{p_1}, \dots, \subst{\theta_n}{p_n}} \in \Sigma,
    \]
    \item \emph{विध्यात्मक अनुमानाखाली} बंद असतो;
    म्हणजे $\varphi$ आणि $\varphi
      \lif \psi \in \Sigma$ असतील, तर $\psi \in \Sigma$.
  \end{enumerate}
\end{defn}

मोडल तर्कशास्त्रांसाठी संबंधी चिन्हार्थमीमांसा
वापरायची असेल, तर सर्व मोडल प्रतिमानांत वैध
असलेली सर्व सूत्रे त्यांत समाविष्ट असावीत,
अशीही अट हवी. \Ax{K} आणि~\Dual{}
चे सर्व दृष्टांत निष्पन्न करता येण्याजोगे असतील, आणि
सूत्र~$\varphi$ निष्पन्न करता येण्याजोगे असेल तेव्हा
$\Box \varphi$ देखील तसेच असेल, तर ही अट पूर्ण होते.
या अटी पाळणाऱ्या मोडल तर्कशास्त्राला
\emph{सामान्य} म्हणतात. (असामान्य मोडल तर्कशास्त्रेही
आहेत; पण नेहमीची संबंधी प्रतिमाने त्यांच्यासाठी
पुरेशी नाहीत.)

\begin{defn}
  मोडल तर्कशास्त्र $\Sigma$ पुढील सूत्रे सामावत असेल,
  \begin{align*}
    \tag{\Ax{K}} & \Box(p \lif q) \lif (\Box p \lif \Box q),
    \\
      \tag{\Dual} & \Diamond p \liff \lnot\Box\lnot p
  \end{align*}
  आणि \emph{अनिवार्यीकरणाखाली} बंद असेल, म्हणजे
  $\varphi \in
  \Sigma$ असेल, तर $\Box \varphi \in \Sigma$,
  तर त्याला \emph{सामान्य} म्हणतात.
\end{defn}

सर्वतः-सत्य अभिव्यंजन \emph{जगातील सत्यता}
टिकवण्यासाठी पुरेसे सूक्ष्म असते; पण \Nec{} हा
नियम फक्त \emph{प्रतिमानातील सत्यता} टिकवतो
(आणि म्हणून चौकटीतील किंवा चौकटींच्या वर्गातील
वैधताही टिकवतो), हे लक्षात घ्या.

\begin{prop}\label{nml:prf:nor:prop:rk}
  प्रत्येक सामान्य मोडल तर्कशास्त्र \RK
  या नियमाखाली बंद असते:
  \begin{prooftree}
    \AxiomC{$\varphi_1 \lif (\varphi_2 \lif \cdots (\varphi_{n-1} \lif \varphi_n)\cdots)$}
    \RightLabel{\RK}
    \UnaryInfC{$\Box\varphi_1 \lif (\Box\varphi_2 \lif \cdots (\Box\varphi_{n-1}
      \lif \Box\varphi_n)\cdots).$}
  \end{prooftree}
\end{prop}

\begin{proof}
  $n$ वर विगमनाने सिद्ध करू. $n = 1$ असेल,
  तर हा नियम नेमका \Nec असतो, आणि प्रत्येक
  सामान्य मोडल तर्कशास्त्र \Nec खाली बंद असते.

  आता निष्कर्ष $n-1$ साठी खरा आहे असे समजा;
  तो~$n$ साठी सिद्ध करू.

  पुढील गृहीत धरा:
  \begin{align*}
  & \varphi_1 \lif (\varphi_2 \lif \cdots (\varphi_{n-1} \lif \varphi_n)\cdots) \in \Sigma
  \intertext{विगमन गृहीतकावरून मिळते:}
  & \Box\varphi_1 \lif (\Box\varphi_2 \lif \cdots \Box(\varphi_{n-1} \lif \varphi_n)\cdots)
  \in \Sigma
  \intertext{$\Sigma$ सामान्य मोडल तर्कशास्त्र असल्यामुळे
    त्यात~$\Ax{K}$ चे सर्व दृष्टांत आहेत; विशेषतः:}
  & \Box(\varphi_{n-1} \lif \varphi_n) \lif (\Box\varphi_{n-1} \lif \Box\varphi_n) \in \Sigma
  \intertext{विध्यात्मक अनुमान आणि योग्य सर्वतः-सत्य दृष्टांत
    वापरून मिळते:}
  & \Box\varphi_1 \lif (\Box\varphi_2 \lif \cdots (\Box\varphi_{n-1}
  \lif \Box\varphi_n)\cdots) \in \Sigma.
  \end{align*}
\end{proof}

\begin{prop}\label{nml:prf:nor:prop:notDiamondBot}
  प्रत्येक सामान्य मोडल तर्कशास्त्र $\Sigma$
  हे~$\lnot\Diamond\lfalse$ सामावते.
\end{prop}

\begin{prob}
  \ref{nml:prf:nor:prop:notDiamondBot} सिद्ध करा.
\end{prob}

\begin{prop}
  $\varphi_1$, \dots, $\varphi_n$ ही सूत्रे असू द्या.
  मग $\varphi_1$, \dots, $\varphi_n$ यांचे सर्व दृष्टांत
  सामावणारे लघुतम सामान्य मोडल तर्कशास्त्र $\Sigma$
  अस्तित्वात असते.
\end{prop}

\begin{proof}
  $\varphi_1$, \dots, $\varphi_n$ दिल्यावर
  $\varphi_1$, \dots, $\varphi_n$ यांचे सर्व दृष्टांत सामावणाऱ्या
  सर्व सामान्य मोडल तर्कशास्त्रांचा
  छेद~$\Sigma$ म्हणून ठरवा. अशा तर्कशास्त्रांचा
  वर्ग रिकामा नाही: सर्व सूत्रांचा संच
  $\Frm[L]$ हेही असे सामान्य मोडल तर्कशास्त्र आहे.
\end{proof}

\begin{defn}
$\varphi_1$, \dots, $\varphi_n$ सामावणाऱ्या लघुतम
सामान्य मोडल तर्कशास्त्राला \emph{मोडल प्रणाली}
म्हणतात; ते $\Log{K} \varphi_1 \dots
\varphi_n$ असे दर्शवतात. लघुतम सामान्य मोडल
तर्कशास्त्र~\Log{K} असे दर्शवतात.
\end{defn}












\section{निष्पत्ती आणि मोडल प्रणाली}\label{nml:prf:prf:sec}

सामान्य मोडल तर्कशास्त्रांसाठी निष्पत्ती
म्हणजे काय ते आधी ठरवू. साधारणपणे,
निष्पत्ती म्हणजे सूत्रांची अशी मालिका
की तिच्यातील प्रत्येक घटक काही
\emph{स्वयंसिद्धकांपैकी} एक (किंवा त्याचा आदेशन-दृष्टांत)
असतो, किंवा आधीच्या घटक पासून मोजक्या
अनुमान नियमांपैकी एखाद्याने मिळतो. सामान्य मोडल
तर्कशास्त्रात सर्व सर्वतः-सत्य सूत्रांचे दृष्टांत
, \Ax{K}, आणि~\Dual{}
स्वयंसिद्धके मानली जातात. त्यातून लघुतम सामान्य
मोडल तर्कशास्त्र, म्हणजे मोडल प्रणाली~$\Log{K}$,
मिळते. इतर प्रणाली मिळवण्यासाठी आणखी स्वयंसिद्धके
जोडता येतात. नियम मात्र नेहमी विध्यात्मक अनुमान~\MP{}
आणि अनिवार्यीकरण~\Nec हेच असतात.

\begin{defn}
  मोडल प्रणाली $\Log{K} \varphi_1 \dots \varphi_n$ आणि
  सूत्र~$\psi$ दिले असता, $\chi_1$, \dots,
  $\chi_k$ अशी सूत्रे असतील, जिथे $\chi_k = \psi$
  आणि प्रत्येक $\chi_i$ हा सर्वतः-सत्य दृष्टांत,
  किंवा $\Ax{K}$, $\Dual$,
  $\varphi_1$, \dots, $\varphi_n$ यांपैकी एखाद्याचा दृष्टांत,
  किंवा आधीच्या सूत्रांपासून \MP{} अथवा~\Nec
  नियमाने मिळणारे सूत्र असेल, तर आणि तेव्हाच
  $\psi$ हे $\Log{K} \varphi_1 \dots
  \varphi_n$ मध्ये \emph{निष्पन्न करता येण्याजोगे} आहे असे म्हणतो;
  संकेतात $\Log{K} \varphi_1 \dots \varphi_n \Proves \psi$.
\end{defn}

पुढील प्रतिज्ञेनुसार $\psi$ ची $\Sigma$-निष्पत्ती
दाखवून $\psi \in \Sigma$ हे सिद्ध करता येते.

\begin{prop}
  $\Log{K} \varphi_1 \dots \varphi_n = \Setabs{\psi}{\Log{K} \varphi_1 \dots \varphi_n \Proves \psi}$.
\end{prop}

\begin{proof}
  निष्पत्त्यांच्या लांबीवर विगमन करून
  $\Setabs{\psi}{\Log{K} \varphi_1 \dots \varphi_n \Proves \psi} \subseteq
  \Log{K} \varphi_1 \dots \varphi_n$ हे दाखवू.

  $\psi$ च्या निष्पत्तीची लांबी~$1$ असेल,
  तर त्यात एकच सूत्र असते. ते सूत्र आधीच्या
  सूत्रांपासून \MP{} किंवा \Nec ने मिळालेले
  असू शकत नाही. म्हणून ते सर्वतः-सत्य दृष्टांत,
  \Ax{K} चा, $\Dual$ चा,
  किंवा $\varphi_1$, \dots,~$\varphi_n$ यांपैकी एखाद्याचा
  दृष्टांत असला पाहिजे. $\Log{K}\varphi_1\dots\varphi_n$
  हे सर्व सामावते; म्हणून $\psi \in \Log{K}\varphi_1
  \dots \varphi_n$.

  $\psi$ च्या निष्पत्तीची लांबी~$> 1$
  असेल, तर आधीच्या प्रकरणांखेरीज निष्पत्तीमधील अंतिम ओळ नसलेल्या सूत्रांपासून
  \MP{} किंवा \Nec{} नेही
  $\psi$ मिळू शकते. $\psi$ हे $\chi$ आणि
  $\chi \lif \psi$ यांपासून (\MP ने) मिळत असेल,
  तर विगमन गृहीतकावरून $\chi$ आणि
  $\chi \lif \psi \in
  \Log{K}\varphi_1 \dots \varphi_n$. प्रत्येक मोडल
  तर्कशास्त्र विध्यात्मक अनुमानाखाली बंद असते;
  म्हणून $\psi \in \Log{K}\varphi_1 \dots
  \varphi_n$. $\psi \equiv \Box \chi$ हे $\chi$
  पासून \Nec ने मिळत असेल, तर विगमन
  गृहीतकावरून $\chi
  \in \Log{K}\varphi_1 \dots \varphi_n$. प्रत्येक
  सामान्य मोडल तर्कशास्त्र $\Nec$ खाली बंद असते;
  म्हणून $\psi \in
  \Log{K}\varphi_1\dots\varphi_n$.

  उलट समावेशासाठी $\Sigma = \Setabs{\psi}{\Log{K} \varphi_1 \dots \varphi_n \Proves \psi }$
  हे $\varphi_1$, \dots, $\varphi_n$ चे सर्व दृष्टांत
  सामावणारे सामान्य मोडल तर्कशास्त्र आहे, हे
  दाखवू. मग $\Log{K} \varphi_1 \dots \varphi_n$
  हे व्याख्येने असे लघुतम तर्कशास्त्र आहे,
  या तथ्याचा उपयोग करू.
  \begin{enumerate}
    \item प्रत्येक सर्वतः-सत्य सूत्र~$\psi$ हाही
      सर्वतः-सत्य दृष्टांत आहे. म्हणून
      $\Log{K}\varphi_1\dots\varphi_n \Proves \psi$;
      त्यामुळे $\Sigma$ सर्व सर्वतः-सत्य सूत्रे सामावते.
    \item $\Log{K}\varphi_1\dots\varphi_n \Proves \chi$
      आणि $\Log{K}\varphi_1\dots\varphi_n \Proves \chi \lif \psi$
      असल्यास $\Log{K}\varphi_1\dots\varphi_n \Proves \psi$:
      $\chi$ ची निष्पत्ती आणि $\chi \lif \psi$
      ची निष्पत्ती जोडून अंतिम ओळ~$\psi$
      घाला. या ओळीला \MP{} चे समर्थन मिळते.
      म्हणून $\Sigma$ विध्यात्मक अनुमानाखाली बंद आहे.
    \item $\psi$ ची निष्पत्ती असेल, तर
      $\psi$ च्या प्रत्येक आदेशन-दृष्टांताचीही
      निष्पत्ती असते: संपूर्ण निष्पत्तीमधील प्रत्येक सूत्र वर तेच आदेशन लावा.
      (सराव: निष्पत्त्यांच्या लांबीवर विगमन
      करून परिणामी मालिकाही योग्य निष्पत्ती
      आहे हे सिद्ध करा.) म्हणून $\Sigma$ एकरूप
      आदेशनाखाली बंद आहे. ($\Sigma$ मोडल तर्कशास्त्राच्या
      सर्व अटी पूर्ण करते, हे आता सिद्ध झाले.)
    \item $\Log{K}\varphi_1\dots\varphi_n \Proves \Ax{K}$;
      म्हणून $\Ax{K} \in \Sigma$.
    \item $\Log{K}\varphi_1\dots\varphi_n \Proves \Dual$
      असल्यामुळे $\Dual \in \Sigma$.
    \item $\Log{K}\varphi_1\dots\varphi_n \Proves \chi$
      असेल, तर अतिरिक्त ओळ~$\Box \chi$ ला~\Nec
      चे समर्थन मिळते. त्यामुळे $\Sigma$ हा
      \Nec खाली बंद आहे; म्हणजे $\Sigma$ सामान्य आहे.
  \end{enumerate}
\end{proof}













\section{\Log{K} मधील सिद्धता}\label{nml:prf:prk:sec}

लघुतम मोडल प्रणालीतील सिद्धतेचा सराव म्हणून
\ref{nml:syn:sch:tab:valid-invalidSchemas} च्या
डाव्या बाजूची वैध सूत्रे सर्व
\Log{K} मध्ये सिद्ध करता येतात, हे दाखवू.

\begin{prop}
  $\Log{K} \Proves \Box\varphi \lif \Box (\psi\lif \varphi)$
\end{prop}

\begin{proof}
  \begin{derivation}
    1. & $\varphi \lif (\psi \lif \varphi)$ & \Taut \\
    2. & $\Box(\varphi \lif (\psi \lif \varphi))$ & \Nec, 1 \\
    3. & $\Box(\varphi \lif (\psi \lif \varphi)) \lif
    (\Box\varphi \lif \Box(\psi \lif \varphi))$ & \Ax{K}\\
    4. & $\Box\varphi \lif \Box (\psi\lif \varphi)$ & \MP, 2, 3
  \end{derivation}
\end{proof}

\begin{prop}
  $\Log{K} \Proves \Box(\varphi \land \psi) \lif (\Box \varphi \land \Box\psi)$
\end{prop}

\begin{proof}
\begin{derivation}
  1. & $(\varphi \land \psi) \lif \varphi$ & \Taut \\
  2. & $\Box((\varphi \land \psi) \lif \varphi)$ & \Nec \\
  3. & $\Box((\varphi \land \psi) \lif \varphi) \lif (\Box(\varphi \land \psi) \lif \Box\varphi)$
  & \Ax{K} \\
  4. & $\Box(\varphi \land \psi) \lif \Box\varphi$ & \MP, 2, 3 \\
  5. & $(\varphi \land \psi) \lif \psi$ & \Taut \\
  6. & $\Box((\varphi \land \psi) \lif \psi)$ & \Nec \\
  7. & $\Box((\varphi \land \psi) \lif \psi) \lif (\Box(\varphi \land \psi) \lif \Box\psi)$
  & \Ax{K} \\
  8. & $\Box(\varphi \land \psi) \lif \Box\psi$ & \MP, 6, 7 \\
  9. & $(\Box(\varphi \land \psi) \lif \Box\varphi) \lif{}$ \\
  & \qquad $((\Box(\varphi \land \psi) \lif \Box\psi) \lif{}$ \\
  & \qquad $(\Box(\varphi \land \psi) \lif (\Box \varphi \land \Box\psi)))$ & \Taut\\
  10. &  $(\Box(\varphi \land \psi) \lif \Box\psi) \lif{}$ \\
  & \qquad $(\Box(\varphi \land \psi) \lif (\Box \varphi \land \Box\psi))$ & \MP, 4, 9\\
  11. & $\Box(\varphi \land \psi) \lif (\Box \varphi \land \Box\psi)$ & \MP, 8, 10.
\end{derivation}
ओळ~$9$ वरील सूत्र पुढील सर्वतः-सत्य सूत्राचा दृष्टांत आहे:
\[(p \lif q) \lif ((p \lif r) \lif (p \lif (q \land r))).
\]
\end{proof}

\begin{prop}
  $\Log{K}\Proves (\Box\varphi \land \Box\psi) \lif \Box (\varphi \land \psi)$
\end{prop}

\begin{proof}
  \begin{derivation}
    1. & $\varphi \lif (\psi \lif (\varphi \land \psi))$ & \Taut \\
    2. & $\Box(\varphi \lif (\psi \lif (\varphi \land \psi)))$ & \Nec, 1 \\
    3. & $\Box(\varphi \lif (\psi \lif (\varphi \land \psi))) \lif (\Box\varphi \lif \Box(\psi \lif (\varphi \land \psi)))$ & \Ax{K}\\
    4. & $\Box\varphi \lif \Box(\psi \lif (\varphi \land \psi))$ & \MP, 2, 3\\
    5. & $\Box(\psi \lif (\varphi \land \psi)) \lif (\Box\psi \lif \Box(\varphi \land \psi))$ & \Ax{K} \\
    6. & $(\Box\varphi \lif \Box(\psi \lif (\varphi \land \psi))) \lif {}$ \\
    & \qquad $((\Box(\psi \lif (\varphi \land \psi)) \lif (\Box\psi \lif \Box(\varphi \land \psi))) \lif {}$\\
    & \qquad $(\Box\varphi \lif (\Box \psi \lif \Box(\varphi \land \psi))))$ & \Taut\\
    7. & $(\Box(\psi \lif (\varphi \land \psi)) \lif (\Box\psi \lif \Box(\varphi \land \psi))) \lif {}$\\
    & \qquad $(\Box\varphi \lif (\Box \psi \lif \Box(\varphi \land \psi)))$ & \MP, 4, 6\\
    8. & $\Box\varphi \lif (\Box \psi \lif \Box(\varphi \land \psi))$ & \MP, 5, 7\\
    9. & $(\Box\varphi \lif (\Box \psi \lif \Box(\varphi \land \psi))) \lif {}$\\
    & \qquad $((\Box\varphi \land \Box\psi) \lif \Box(\varphi \land \psi))$ & \Taut\\
    10. & $(\Box\varphi \land \Box\psi) \lif \Box(\varphi \land \psi)$ & \MP, 8, 9
  \end{derivation}
  ओळी $6$ आणि $9$ वरील सूत्रे अनुक्रमे पुढील
  सर्वतः-सत्य सूत्रांचे दृष्टांत आहेत:
  \begin{align*}
    (p \lif q) & \lif ((q \lif r) \lif (p \lif r)) \\
    (p \lif (q \lif r)) & \lif ((p \land q) \lif r)
  \end{align*}
\end{proof}

\begin{prop}
  $\Log{K} \Proves \lnot\Box p \lif \Diamond \lnot
    p$
\end{prop}

\begin{proof}

  \begin{derivation}
    1. & $\Diamond \lnot p \liff \lnot\Box\lnot\lnot p$ & \Dual\\
    2. & $(\Diamond \lnot p \liff \lnot\Box\lnot\lnot p) \lif {}$ \\
    & \qquad $(\lnot \Box \lnot\lnot p \lif \Diamond \lnot p)$& \Taut\\
    3. & $\lnot \Box \lnot\lnot p \lif \Diamond \lnot p$ & \MP, 1, 2\\
    4. & $\lnot\lnot p \lif p$ & \Taut\\
    5. & $\Box(\lnot\lnot p \lif p)$ & \Nec, 4\\
    6. & $\Box(\lnot\lnot p \lif p) \lif (\Box \lnot\lnot p \lif \Box p)$ & \Ax{K}\\
    7. & $(\Box \lnot\lnot p \lif \Box p)$ & \MP, 5, 6\\
    8. & $(\Box \lnot\lnot p \lif \Box p) \lif (\lnot \Box p \lif \lnot\Box\lnot\lnot p)$ & \Taut\\
    9. & $\lnot \Box p \lif \lnot\Box\lnot\lnot p$ & \MP, 7, 8\\
    10. & $(\lnot \Box p \lif \lnot\Box\lnot\lnot p) \lif {}$ \\
    & \qquad $((\lnot \Box \lnot\lnot p \lif \Diamond \lnot p) \lif (\lnot \Box p \lif \Diamond\lnot p))$ & \Taut\\
    11. & $(\lnot \Box \lnot\lnot p \lif \Diamond \lnot p) \lif (\lnot \Box p \lif \Diamond\lnot p)$ & \MP, 9, 10\\
    12. & $\lnot\Box p \lif \Diamond \lnot p$ & \MP, 3, 11\\
  \end{derivation}
  ओळी $8$ आणि $10$ वरील सूत्रे अनुक्रमे पुढील
  सर्वतः-सत्य सूत्रांचे दृष्टांत आहेत:
  \begin{align*}
    & (p \lif q) \lif (\lnot q \lif \lnot p) \\
    & (p \lif q) \lif ((q \lif r) \lif (p \lif r)).
  \end{align*}

\end{proof}

\begin{prob}
  पुढील सूत्रे साठी $\Log{K}$ मध्ये
  निष्पत्त्या शोधा:
  \begin{enumerate}
    \item $\Box \lnot p \lif \Box(p \lif q)$
    \item $(\Box p \lor \Box q) \lif \Box(p \lor q)$
    \item $\Diamond p \lif \Diamond(p \lor q)$
  \end{enumerate}
\end{prob}













\section{व्युत्पन्न नियम}\label{nml:prf:der:sec}

निष्पत्त्या शोधणे आणि लिहिणे कठीण, किचकट
आणि पुनरावृत्तीचे काम आहे. उदाहरणार्थ, $\varphi \lif \psi$
पासून $\Box \varphi \lif \Box \psi$ कडे जायचे असेल,
म्हणजे नियम~\RK लावायचा असेल, तर आधी \Nec लावावा लागतो,
मग~\Ax{K} चा योग्य दृष्टांत लिहावा लागतो आणि नंतर~\MP
लावावा लागतो. तसेच $\varphi \lif \psi$ आणि $\psi \lif \chi$
पासून $\varphi \lif \chi$ कडे जाण्यासाठी पुढील (लांब)
सर्वतः-सत्य दृष्टांत लिहावा लागतो:
\[(\varphi \lif \psi) \lif ((\psi \lif \chi) \lif (\varphi \lif \chi))
\]
आणि \MP{} दोनदा लावावा लागतो. अनेकदा एखादे उप-सूत्र
त्याच्याशी समतुल्य असलेल्या सूत्राने बदलायचे असते;
उदाहरणार्थ, $\Diamond
  \varphi$ च्या जागी $\lnot\Box\lnot \varphi$, किंवा
$\lnot\lnot \varphi$ च्या जागी $\varphi$. म्हणून संपूर्ण
निष्पत्ती लिहिण्याऐवजी मध्यवर्ती पायऱ्या
निष्पन्न करता येण्याजोगे का आहेत, एवढेच नोंदवणे सोयीचे आहे.
त्यासाठी निष्पन्नता विषयी काही तथ्ये गोळा करू.

\begin{prop}
  $\Log{K} \Proves \varphi_1$, \dots, $\Log{K} \Proves \varphi_n$
  असेल आणि विधानीय तर्कशास्त्रानुसार $\varphi_1$, \dots,
  $\varphi_n$ पासून $\psi$ निष्पन्न होत असेल, तर
  $\Log{K} \Proves \psi$.
\end{prop}

\begin{proof}
  विधानीय तर्कशास्त्रानुसार $\varphi_1$, \dots,~$\varphi_n$
  पासून $\psi$ निष्पन्न होत असेल, तर
  \[  \varphi_1 \lif (\varphi_2 \lif \cdots (\varphi_n \lif \psi)\dots)
  \]
  हा सर्वतः-सत्य दृष्टांत असतो. \MP{} $n$ वेळा
  लावल्याने $\psi$ ची निष्पत्ती मिळते.
\end{proof}

या प्रतिज्ञेचा वापर~\PL ने दर्शवू.

\begin{prop}
  $\Log{K} \Proves \varphi_1 \lif (\varphi_2 \lif \cdots (\varphi_{n-1} \lif
  \varphi_n)\dots)$ असेल, तर $\Log{K} \Proves \Box \varphi_1 \lif (\Box \varphi_2 \lif
  \cdots (\Box \varphi_{n-1} \lif \Box \varphi_n)\dots)$.
\end{prop}

\begin{proof}
  \ref{nml:prf:nor:prop:rk} च्या सिद्धतेप्रमाणेच
  $n$ वर विगमन करा.
\end{proof}

या प्रतिज्ञेचा वापर~\RK ने दर्शवू. या परिणामांमुळे
निष्पन्नता विषयीची विधाने किती सोप्या
रीतीने सिद्ध करता येतात, ते पाहू.

\begin{prop}
  $\Log{K} \Proves (\Box\varphi \land \Box\psi) \lif \Box (\varphi \land \psi)$
\end{prop}

\begin{proof}
  \begin{derivation}
    1. & $\Log{K} \Proves \varphi \lif (\psi \lif (\varphi \land \psi))$ & \Taut \\
    2. & $\Log{K} \Proves \Box\varphi \lif (\Box \psi \lif \Box(\varphi \land \psi))$ & \RK, 1\\
    3. & $\Log{K} \Proves (\Box\varphi \land \Box\psi) \lif \Box(\varphi \land \psi)$ & \PL, 2
  \end{derivation}
\end{proof}

\begin{prop}\label{nml:prf:der:prop:rewriting}
  $\Log{K} \Proves \varphi \liff \psi$ आणि $\Log{K} \Proves
  \Subst{\chi}{\varphi}{q}$ असेल, तर $\Log{K} \Proves \Subst{\chi}{\psi}{q}$.
\end{prop}

\begin{proof}
  सराव.
\end{proof}

\begin{prob}
  $\chi$ च्या जटिलतेवर विगमन करून
  \ref{nml:prf:der:prop:rewriting} सिद्ध करा:
  $\Log{K} \Proves \varphi \liff \psi$ असेल, तर
  $\Log{K} \Proves \Subst{\chi}{\varphi}{q} \liff \Subst{\chi}{\psi}{q}$
  हे दाखवा.
\end{prob}

$\Diamond$ चे $\Box$ मध्ये (किंवा उलट) रूपांतर करायचे
असेल, किंवा सूत्र मधील दुहेरी निषेध काढायचा
असेल, तेव्हा ही प्रतिज्ञा विशेष उपयोगी ठरते.
पुढे $\chi(\varphi)$ च्या जागी सूत्र~$\chi(\psi)$ लिहिताना
\ref{nml:prf:der:prop:rewriting} चा वापर ``$\psi$ for $\varphi$''
असा दाखवू. दुसऱ्या शब्दांत, ``$\psi$ for $\varphi$''
याचा संक्षिप्त अर्थ असा:
  \begin{derivation}
    & $\Proves \chi(\varphi)$ \\
    & $\Proves \varphi \liff \psi$\\
    & $\Proves \chi(\psi)$ & \ref{nml:prf:der:prop:rewriting} वरून
  \end{derivation}
उदाहरणार्थ:

\begin{prop}
  $\Log{K} \Proves \lnot\Box p \lif \Diamond \lnot p$
\end{prop}

\begin{proof}

  \begin{derivation}
    1. & $\Log{K} \Proves \Diamond \lnot p \liff \lnot\Box\lnot\lnot p$ &
    \Dual\\
    2. & $\Log{K} \Proves \lnot \Box \lnot\lnot p \lif \Diamond \lnot p$ & \PL, 1\\
    3. & $\Log{K} \Proves \lnot\Box p \lif \Diamond \lnot p$ & $p$ for $\lnot\lnot p$\\
  \end{derivation}

\end{proof}

वरील निष्पत्तीमधील अंतिम पायरी
``$p$ for $\lnot\lnot p$'' ही पुढील पायऱ्यांचे
संक्षिप्त रूप आहे:
  \begin{derivation}
    & $\Log{K} \Proves \lnot \Box \lnot\lnot p \lif \Diamond \lnot
    p$\\
    & $\Log{K} \Proves \lnot\lnot p \liff p$ & \Taut\\
    & $\Log{K} \Proves \lnot\Box p \lif \Diamond \lnot p$ & \ref{nml:prf:der:prop:rewriting} वरून
  \end{derivation}
\ref{nml:prf:der:prop:rewriting} मधील $\chi(q)$, $\varphi$ आणि $\psi$
यांच्या जागी येथे अनुक्रमे
$\lnot \Box q \lif \Diamond \lnot p$,
$\lnot\lnot p$ आणि~$p$ आहेत.

सूत्रामध्ये $\lnot\Diamond \varphi$ हे उप-सूत्र
असेल, तर \ref{nml:prf:der:prop:rewriting} वापरून त्याच्या जागी
$\Box\lnot \varphi$ लिहिता येते, कारण
$\Log{K} \Proves \lnot\Diamond \varphi \liff \Box\lnot \varphi$.
हा आणि यांसारखे बदल थोडक्यात
``$\Box\lnot$ for $\lnot\Diamond$'' असे दर्शवू.

पुढील प्रतिज्ञेमुळे निष्पन्नता विषयीचे परिणाम
रूपबंधात्मक पद्धतीने सिद्ध करता येतात. उदाहरणार्थ,
आधीची प्रतिज्ञा केवळ $\Log{K} \Proves \lnot\Box p \lif
\Diamond \lnot p$ एवढेच दाखवत नाही; कोणत्याही~$\varphi$
साठी $\Log{K} \Proves \lnot\Box \varphi \lif \Diamond
\lnot \varphi$ हेही दाखवते.

\begin{prop}
  $\varphi$ हा $\psi$ चा आदेशन-दृष्टांत असेल आणि
  $\Log{K} \Proves \psi$ असेल, तर $\Log{K} \Proves \varphi$.
\end{prop}

\begin{proof}
  संपूर्ण निष्पत्ती वर आदेशन लावल्यानेही योग्य
  निष्पत्ती मिळते, हे तपासणे किचकट पण नेहमीचे
  काम आहे ($\psi$ च्या निष्पत्तीच्या लांबीवर
  विगमन करा). विशेषतः, सर्वतः-सत्य दृष्टांतांचे
  आदेशन-दृष्टांतही सर्वतः-सत्य दृष्टांत असतात;
  \Dual{} आणि~\Ax{K}
  यांच्या दृष्टांतांचे आदेशन-दृष्टांतही
  \Dual{} आणि~\Ax{K}
  यांचेच दृष्टांत असतात; आणि आधारसूत्रांमध्ये
  व निष्कर्षात विधानीय चलांच्या
  जागी सूत्रे ठेवले तरी \MP{} व~\Nec{}
  यांचा वापर योग्यच राहतो.
\end{proof}













\section{\Log{K} मधील आणखी सिद्धता}\label{nml:prf:mpr:sec}

\ref{nml:prf:der:sec} मध्ये दिलेली सुलभ पद्धत वापरून
\Log{K} मधील निष्पन्नता ची आणखी काही उदाहरणे पाहू.

\begin{prop}
  $\Log{K} \Proves \Box (\varphi \lif \psi) \lif (\Diamond \varphi \lif \Diamond
  \psi)$
\end{prop}

\begin{proof}
\begin{derivation}
1. & $\Log{K} \Proves (\varphi \lif \psi) \lif (\lnot \psi \lif \lnot \varphi)$ & \PL \\
2. &  $\Log{K} \Proves \Box(\varphi \lif \psi) \lif (\Box\lnot\psi \lif \Box\lnot\varphi)$ & \RK, 1 \\
3. & $\Log{K} \Proves (\Box\lnot\psi \lif \Box\lnot\varphi) \lif (\lnot
\Box\lnot\varphi \lif \lnot\Box\lnot\psi)$ & \Taut\\
4. & $\Log{K} \Proves \Box(\varphi \lif \psi) \lif (\lnot
\Box\lnot\varphi \lif \lnot\Box\lnot\psi)$ & \PL, 2, 3 \\
5. &  $\Log{K} \Proves \Box(\varphi \lif \psi) \lif (\Diamond\varphi \lif \Diamond\psi)$ & $\Diamond$ for $\lnot\Box\lnot$.
\end{derivation}
\end{proof}

\begin{prop}
$\Log{K}\Proves \Box\varphi \lif (\Diamond(\varphi \lif \psi) \lif
  \Diamond \psi)$
\end{prop}

\begin{proof}
  \begin{derivation}
  1. & $\Log{K} \Proves \varphi \lif (\lnot\psi \lif \lnot (\varphi \lif \psi))$ & \Taut \\
  2. & $\Log{K} \Proves \Box\varphi \lif  (\Box\lnot\psi \lif \Box\lnot (\varphi\lif \psi))$ &
  \RK, 1 \\
  3. &  $\Log{K} \Proves \Box\varphi \lif (\lnot \Box\lnot (\varphi\lif \psi) \lif \lnot
  \Box\lnot\psi)$ & \PL, 2 \\
  4. & $\Log{K} \Proves \Box\varphi \lif (\Diamond(\varphi \lif \psi) \lif
  \Diamond \psi)$ & $\Diamond$ for $\lnot\Box\lnot$.
  \end{derivation}
\end{proof}

\begin{prop}
  $\Log{K} \Proves (\Diamond\varphi \lor \Diamond\psi) \lif \Diamond(\varphi \lor \psi)$
\end{prop}

\begin{proof}
  \begin{derivation}
    1. & $\Log{K} \Proves \lnot(\varphi \lor \psi) \lif \lnot\varphi$ & \Taut \\
    2. & $\Log{K} \Proves \Box\lnot(\varphi \lor \psi) \lif \Box\lnot\varphi$ & \RK, 1 \\
    3. & $\Log{K} \Proves \lnot\Box\lnot \varphi \lif \lnot\Box\lnot(\varphi \lor\psi)$ &
    \PL, 2\\
    4. & $\Log{K} \Proves \Diamond\varphi \lif \Diamond(\varphi \lor \psi)$ &
    $\Diamond$ for $\lnot\Box\lnot$\\
    5. & $\Log{K} \Proves \Diamond\psi \lif \Diamond(\varphi \lor \psi)$ & तसेच\\
    6. & $\Log{K} \Proves (\Diamond\varphi \lor\Diamond\psi) \lif \Diamond(\varphi \lor \psi)$
    & \PL, 4, 5.
  \end{derivation}
\end{proof}

\begin{prop}
  $\Log{K} \Proves \Diamond(\varphi \lor\psi) \lif (\Diamond\varphi \lor \Diamond\psi)$
\end{prop}

\begin{proof}
  \begin{derivation}
    1. & $\Log{K} \Proves \lnot \varphi \lif (\lnot \psi \lif \lnot (\varphi \lor \psi))$ & \Taut \\
    2. & $\Log{K} \Proves \Box\lnot \varphi \lif
    (\Box\lnot\psi \lif \Box \lnot (\varphi \lor \psi))$ & \RK\\
    3. & $\Log{K} \Proves \Box\lnot \varphi \lif (\lnot \Box \lnot (\varphi \lor\psi)
    \lif \lnot\Box\lnot \psi)$ & \PL, 2\\
    4. & $\Log{K} \Proves \lnot \Box \lnot(\varphi \lor \psi) \lif (\Box \lnot \varphi \lif
    \lnot\Box\lnot \psi)$ & \PL, 3\\
    5. & $\Log{K} \Proves \lnot \Box \lnot(\varphi\lor\psi) \lif (\lnot
    \lnot\Box\lnot\psi \lif \lnot\Box\lnot\varphi)$ & \PL, 4\\
    6. & $\Log{K} \Proves \Diamond(\varphi \lor \psi) \lif (\lnot
    \Diamond\psi \lif \Diamond\varphi)$ & $\Diamond$ for $\lnot\Box\lnot$\\
    7. & $\Log{K} \Proves \Diamond(\varphi\lor\psi) \lif (\Diamond\psi \lor \Diamond\varphi)$ & \PL, 6. \\
  \end{derivation}
\end{proof}

\begin{prob}
  पुढील निष्पन्नता विषयीची विधाने सिद्ध करा:
  \begin{enumerate}
  \item $\Log{K} \Proves \Diamond \lnot \lfalse \lif (\Box \varphi \lif
    \Diamond \varphi)$;
  \item $\Log{K} \Proves \Box(\varphi \lor \psi) \lif (\Diamond \varphi \lor \Box
    \psi)$;
  \item $\Log{K} \Proves (\Diamond \varphi \lif \Box \psi) \lif \Box(\varphi \lif
    \psi)$.
  \end{enumerate}
\end{prob}












\section{द्वयी सूत्र}\label{nml:prf:dua:sec}

\begin{defn}\label{nml:prf:dua:def:duals}
  \Ax{T}, \Ax{B}, \Ax{4} आणि \Ax{5} ही
  सूत्रे आहेत; त्यांपैकी प्रत्येकाला
  \emph{द्वयी} सूत्र आहे.
  ते तळसूचक हिऱ्याने पुढीलप्रमाणे दर्शवले जाते:
  \begin{align}
    \tag{\Ax{T_\Diamond}} p & \lif \Diamond p\\
    \tag{\Ax{B_\Diamond}} \Diamond\Box p & \lif p\\
    \tag{\Ax{4_\Diamond}} \Diamond\Diamond p & \lif \Diamond p\\
    \tag{\Ax{5_\Diamond}} \Diamond\Box p & \lif \Box p
    \end{align}
\end{defn}

वरील द्वयी सूत्रांपैकी प्रत्येक सूत्र संबंधित
सूत्र मधील $p$ च्या जागी $\lnot p$ ठेवून,
प्रतिपरिवर्तन करून,
$\lnot\Box\lnot$ च्या जागी $\Diamond$ आणि
$\lnot\Diamond\lnot$ च्या जागी~$\Box$ ठेवून मिळते.
\Ax{D}, म्हणजे $\Box\varphi \lif \Diamond\varphi$, हे या अर्थाने
स्वतःचेच द्वयी सूत्र आहे.

\begin{prop}\label{nml:prf:dua:prop:dualsys}
  \ref{nml:prf:dua:def:duals} मधील प्रत्येक
  सूत्र~$\varphi$ साठी:
  $\Log{K}\varphi = \Log{K}\varphi_{\Diamond}$.
\end{prop}
\begin{proof}
  सराव.
\end{proof}
\begin{prob}
  \ref{nml:prf:dua:prop:dualsys} सिद्ध करा.
\end{prob}












\section{मोडल प्रणालींमधील सिद्धता}\label{nml:prf:prs:sec}

आता~\Log{K} व्यतिरिक्त इतर मोडल तर्कशास्त्रीय
प्रणालींमधील सिद्धता पाहू.

\begin{prop}\label{nml:prf:prs:prop:S5facts}
पुढील सिद्धताविषयक परिणाम मिळतात:
  \begin{enumerate}
  \item $\Log{KT5} \Proves \Ax{B}$;
  \item $\Log{KT5} \Proves \Ax{4}$;
  \item $\Log{KDB4} \Proves \Ax{T}$;
  \item $\Log{KB4} \Proves \Ax{5}$;
  \item $\Log{KB5} \Proves \Ax{4}$;
  \item \label{nml:prf:prs:prop:S5facts-KT-D}$\Log{KT} \Proves \Ax{D}$.
  \end{enumerate}
\end{prop}

\begin{proof}
  प्रत्येक परिणामाची सिद्धता देऊ.
  \begin{enumerate}
    \item $\Log{KT5} \Proves \Ax{B}$:
      \begin{derivation}
        1. & $\Log{KT5} \Proves \Diamond\varphi \lif \Box\Diamond\varphi$ & \Ax{5}\\
        2. & $\Log{KT5} \Proves \varphi \lif \Diamond\varphi$ & $\Ax{T_\Diamond}$\\
        3. & $\Log{KT5} \Proves \varphi \lif \Box\Diamond\varphi$ & \PL, 2, 1.
      \end{derivation}
    \item $\Log{KT5} \Proves \Ax{4}$:
      \begin{derivation}
        1. &$\Log{KT5} \Proves \Diamond\Box\varphi \lif \Box\Diamond\Box\varphi$ & \Ax{5}
        मध्ये $p$ च्या जागी $\Box\varphi$\\
        2. & $\Log{KT5} \Proves \Box\varphi \lif \Diamond\Box\varphi$ & $\Ax{T_\Diamond}$
        मध्ये $p$ च्या जागी $\Box\varphi$\\
        3. & $\Log{KT5} \Proves \Box\varphi \lif \Box\Diamond\Box\varphi$ & \PL, 2, 1\\
        4. & $\Log{KT5} \Proves \Diamond\Box\varphi \lif \Box\varphi$ & $\Ax{5_\Diamond}$\\
        5. & $\Log{KT5} \Proves \Box\Diamond\Box\varphi \lif \Box\Box\varphi$ & \RK{}, 4 \\
        6. & $\Log{KT5} \Proves \Box\varphi \lif \Box\Box\varphi$ & \PL, 3, 5. \\
      \end{derivation}
    \item $\Log{KDB4} \Proves \Ax{T}$:
      \begin{derivation}
        1. & $\Log{KDB4} \Proves \Diamond\Box\varphi \lif \varphi $ & $\Ax{B_\Diamond}$ \\
        2. & $\Log{KDB4} \Proves \Box\Box\varphi \lif \Diamond\Box\varphi$ & $\Ax{D}$
        मध्ये $p$ च्या जागी $\Box\varphi$\\
        3. & $\Log{KDB4} \Proves \Box\Box\varphi \lif \varphi$ & \PL, 1, 2\\
        4. & $\Log{KDB4} \Proves \Box\varphi \lif \Box\Box\varphi$ & \Ax{4} \\
        5. & $\Log{KDB4} \Proves \Box\varphi \lif \varphi$ & \PL, 4, 3. \\
      \end{derivation}
    \item $\Log{KB4} \Proves \Ax{5}$:
      \begin{derivation}
        1. & $\Log{KB4} \Proves \Diamond\varphi \lif \Box \Diamond\Diamond\varphi$ & \Ax{B}
        मध्ये $p$ च्या जागी $\Diamond\varphi$\\
        2. & $\Log{KB4} \Proves \Diamond\Diamond \varphi \lif \Diamond \varphi$ &
        $\Ax{4_\Diamond}$ \\
        3. & $\Log{KB4} \Proves \Box\Diamond\Diamond \varphi \lif \Box \Diamond\varphi$ &
        \RK{}, 2 \\
        4. & $\Log{KB4} \Proves \Diamond\varphi \lif \Box\Diamond\varphi$ & \PL, 1, 3.
      \end{derivation}
    \item $\Log{KB5} \Proves \Ax{4}$:
      \begin{derivation}
        1. & $\Log{KB5} \Proves \Box\varphi \lif \Box\Diamond\Box \varphi$ & \Ax{B} मध्ये
        $p$ च्या जागी $\Box\varphi$ \\
        2. & $\Log{KB5} \Proves  \Diamond\Box\varphi \lif \Box \varphi$ &
        $\Ax{5_\Diamond}$\\
        3. & $\Log{KB5} \Proves  \Box\Diamond\Box\varphi \lif \Box\Box \varphi$ & \RK{}, 2 \\
        4. & $\Log{KB5} \Proves  \Box\varphi \lif \Box\Box\varphi$ & \PL, 1, 3.
      \end{derivation}
    \item $\Log{KT} \Proves \Ax{D}$:
      \begin{derivation}
        1. & $\Log{KT} \Proves \Box \varphi \lif \varphi$ & \Ax{T} \\
        2. & $\Log{KT} \Proves \varphi \lif \Diamond\varphi$ & $\Ax{T_\Diamond}$ \\
        3. & $\Log{KT} \Proves \Box \varphi \lif \Diamond \varphi$ &  \PL, 1, 2
      \end{derivation}
  \end{enumerate}
\end{proof}

\begin{defn}
  प्रचलित संकेतानुसार \Log{S4} ही प्रणाली
  \Log{KT4} आणि \Log{S5} ही प्रणाली
  \Log{KTB4} अशी परिभाषित करतो.
\end{defn}

पुढील प्रतिज्ञा दाखवते की अभिजात प्रणाली~\Log{S5}
अनेक समतुल्य स्वयंसिद्धक-प्रणालींनी मिळते. हे
आश्चर्यकारक नाही; स्वयंसिद्धकांचे हे वेगवेगळे संयोग
समतुल्यता संबंधांनाच लक्षणित करतात
(\ref{nml:frd:es5:prop:equivalences} पाहा).

\begin{prop}\label{nml:prf:prs:prop:S5}
  $\Log{KTB4} = \Log{KT5} = \Log{KDB4} = \Log{KDB5}$.
\end{prop}

\begin{proof}
  सराव.
\end{proof}

\begin{prob}
  \ref{nml:prf:prs:prop:S5} सिद्ध करा.
\end{prob}














\section{निर्दोषता}\label{nml:prf:snd:sec}

ज्या निष्पत्ती प्रणालीमध्ये निष्पन्न
करता येणारे प्रत्येक सूत्र वैध असते, तिला निर्दोष
म्हणतात. मोडल प्रणाली विचारात घेताना, म्हणजे
\Ax{K} बरोबरच काही सूत्रे $\varphi_1$,
\dots,~$\varphi_n$ यांचे दृष्टांत वापरता येणाऱ्या
निष्पत्त्या विचारात घेताना, $\varphi_1$, \dots,
$\varphi_n$ यांचे सर्व आदेशन-दृष्टांत सर्व जगांत सत्य असलेल्या
प्रत्येक प्रतिमानात प्रत्येक निष्पन्न करता येण्याजोगे सूत्र सत्य असावे,
अशी अपेक्षा आहे.

\begin{thm}[निर्दोषता प्रमेय]\label{nml:prf:snd:thm:soundness}
  $\varphi_1$, \dots, $\varphi_n$ यांचा प्रत्येक दृष्टांत
  अनुक्रमे प्रतिमानांच्या $\mClass{C}_1$, \dots,
  $\mClass{C}_n$ या वर्गांमध्ये वैध असेल,
  तर $\Log{K}\varphi_1\dots \varphi_n
  \Proves \psi$ यावरून $\psi$ हे प्रतिमानांच्या
  $\mClass{C}_1 \cap \dots \cap \mClass{C}_n$
  या वर्गात वैध ठरते.
\end{thm}

\begin{proof}
  सिद्धतांच्या लांबीवर विगमन करू. संक्षेपासाठी
  $\mClass{C} =
  \mClass{C}_1 \cap \dots \cap \mClass{C}_n$ असे ठेवा.
  \begin{enumerate}
  \item विगमनाचा पाया: $\psi$ ची सिद्धता लांबी~$1$
    ची असेल, तर ते सर्वतः-सत्य दृष्टांत,
    \Ax{K} चा, किंवा~\Dual{} चा,
    अथवा $\varphi_1$, \dots,~$\varphi_n$ यांपैकी
    एखाद्याचा दृष्टांत असतो. पहिल्या प्रकरणात
    \ref{nml:syn:tau:prop:valid-taut} वरून $\psi$
    हा $\mClass{C}$ मध्ये वैध आहे, कारण सर्वतः-सत्य
    दृष्टांत प्रतिमानांच्या \emph{कोणत्याही} वर्गात
    वैध असतात. दुसऱ्या प्रकरणात
    \ref{nml:syn:sch:prop:Kvalid} आणि
      \ref{nml:syn:sch:prop:Dual-valid} वरून तेच
    मिळते. शेवटी तिसऱ्या प्रकरणात $\psi$ हा
    $\mClass{C}_i$ मध्ये वैध आहे आणि
    $\mClass{C} \subseteq \mClass{C}_i$ आहे;
    म्हणून \ref{nml:syn:val:prop:subset-class}
    वरून $\psi$ हा $\mClass{C}$ मध्येही वैध आहे.
  \item विगमन पायरी: $\psi$ ची सिद्धता लांबी $k>1$
    ची आहे असे समजा. $\psi$ हा सर्वतः-सत्य दृष्टांत,
    मागील पायरीत नमूद केलेल्या मोडल स्वयंसिद्धकांपैकी
    एखाद्याचा दृष्टांत, किंवा $\varphi_1$, \dots,
    $\varphi_n$ यांपैकी एखाद्याचा दृष्टांत असेल,
    तर मागील पायरीप्रमाणे युक्तिवाद लागू होतो.
    म्हणून आधीच्या
    $\chi \lif \psi$ आणि $\chi$ या सूत्रे वर
    \MP{} लावून $\psi$ मिळाला असे समजा. मग
    $\chi \lif \psi$ आणि $\chi$ यांच्या सिद्धतांची
    लांबी $<k$ आहे; विगमन गृहीतकावरून ती
    $\mClass{C}$ मध्ये वैध आहेत.
    \ref{nml:syn:sch:prop:soundMP} वरून $\psi$
    हाही $\mClass{C}$ मध्ये वैध आहे. शेवटी
    $\chi$ वर \Nec{} लावून $\psi$ मिळाला असे समजा
    (म्हणजे $\psi = \Box\chi$). विगमन गृहीतकावरून
    $\chi$ हा $\mClass{C}$ मध्ये वैध आहे आणि
    \ref{nml:syn:val:prop:Nec-rule} वरून $\psi$ ही
    वैध आहे.
  \end{enumerate}
\end{proof}












\section{प्रणाली भिन्न आहेत हे दाखवणे}\label{nml:prf:dis:sec}

\ref{nml:prf:prs:sec} मध्ये दोन मोडल तर्कशास्त्रीय
प्रणाली प्रत्यक्षात एकच आहेत, हे कसे सिद्ध करायचे
ते पाहिले. आता~\ref{nml:prf:snd:thm:soundness}
वापरून दोन मोडल प्रणाली $\Sigma$ आणि $\Sigma'$
भिन्न आहेत, हे दाखवता येते: असे सूत्र~$\varphi$
शोधायचे की $\Sigma' \Proves \varphi$, पण ते~$\Sigma$
च्या प्रतिमानात असत्य ठरते.

\begin{prop}
  $\Log{KD} \subsetneq \Log{KT}$
\end{prop}

\begin{proof} प्रत्येक परावर्ती संबंध सर्वत्र उत्तरवर्ती असतो,
  या अर्थविषयक तथ्याचा हा वाक्यरचनात्मक समकक्ष आहे.
  $\Log{KD} \subseteq \Log{KT}$ दाखवण्यासाठी
  $\Log{KD} \Proves \psi$ यावरून
  $\Log{KT} \Proves \psi$ मिळते हे पाहावे लागेल;
  \ref{nml:prf:prs:prop:S5facts}\ref{nml:prf:prs:prop:S5facts-KT-D}
  मध्ये दाखवल्याप्रमाणे $\Log{KT} \Proves
  \Ax{D}$ यावरून हे मिळते. हा समावेश कडक आहे हे
  दाखवण्यासाठी निर्दोषतेनुसार
  (\ref{nml:prf:snd:thm:soundness}) \Log{KD} चे असे
  प्रतिमान दाखवणे पुरेसे आहे की त्यात
  \Ax{T}, म्हणजे $\Box p \lif p$, असत्य ठरते
  (हे सोपे काम सरावासाठी ठेवले आहे). मग निर्दोषतेवरून
  $\Log{KD} \Proves/ \Box p \lif p$.
\end{proof}

\begin{prop}
  $\Log{KB} \neq \Log{K4}$.
\end{prop}

\begin{proof}
  \Ax{4} चा एखादा दृष्टांत असत्य ठरणारे सममित
  प्रतिमान रचू. तो दृष्टांत \Log{K4} मध्ये
  निष्पन्न करता येण्याजोगे आहे, पण \Log{KB} मध्ये नाही;
  म्हणून $\Log{K4} \nsubseteq \Log{KB}$ मिळेल.
  \ref{nml:prf:dis:fig:Bnot4} मधील सममित प्रतिमान
  $\mModel{M}$ पाहा. प्रतिमान सममित असल्याने
  \Ax{K} आणि~\Ax{B} ही $\mModel{M}$ मध्ये सत्य
  आहेत (अनुक्रमे \ref{nml:syn:sch:prop:Kvalid}
  आणि \ref{nml:frd:acc:thm:soundschemas} वरून).
  मात्र $\mSat/{M}{\Box p \lif \Box\Box p}[w_1]$.
\end{proof}

\begin{figure}[htpb]
  \centering
  \begin{tikzpicture}[modal]
    \node[world] (w1) [label=above:\mFalse{p},
    label={below:$\mSat{{}}{\Box p}$\\ $\mSat/{{}}{\Box\Box p}$}] {$w_1$};
    \node[world] (w2) [label=above:\mTrue{p},
      label=below:$\mSat/{{}}{\Box p}$, right=of w1]{$w_2$};
    \draw[->, bend left] (w1) to (w2);
    \draw[->, bend left] (w2) to (w1);
  \end{tikzpicture}
  \caption{\Ax{4} चा दृष्टांत असत्य ठरणारे सममित प्रतिमान.}
  \label{nml:prf:dis:fig:Bnot4}
\end{figure}

\begin{thm}\label{nml:prf:dis:thm:KTBnot45}
  $\Log{KTB} \Proves/ \Ax{4}$ आणि $\Log{KTB} \Proves/ \Ax{5}$.
\end{thm}

\begin{proof}
  \ref{nml:frd:acc:thm:soundschemas} वरून
  \Ax{T} आणि \Ax{B} यांचे सर्व दृष्टांत
  अनुक्रमे प्रत्येक परावर्ती आणि सममित प्रतिमानात
  सत्य असतात. म्हणून निर्दोषतेनुसार असे
  परावर्ती सममित प्रतिमान शोधणे पुरेसे आहे की
  त्यातील एका जगात~\Ax{4} चा एखादा दृष्टांत
  असत्य ठरेल; \Ax{5} साठीही तसेच.
  दोन्ही दाव्यांसाठी एकच प्रतिमान वापरू.
  \ref{nml:prf:dis:fig:KTBnot45} मधील सममित, परावर्ती
  प्रतिमान पाहा. येथे $\mSat/{M}{\Box p \to \Box \Box
    p}[w_1]$; म्हणून \Ax{4} हा $w_1$ येथे
  असत्य ठरतो. तसेच $\mSat/{M}{\Diamond
    \lnot p \lif \Box \Diamond \lnot p}[w_2]$;
  म्हणून $\varphi = \lnot p$ असलेला~\Ax{5} चा
  दृष्टांत $w_2$ येथे असत्य ठरतो.
\end{proof}

\begin{figure}[htpb]
  \centering
  \begin{tikzpicture}[modal]
    \node[world] (w1) [label={right:\mTrue{p}},
      label={below: $\mSat{{}}{\Box p}$\\
        $\mSat/{{}}{\Box\Box p}$\\
        $\mSat/{{}}{\Diamond\lnot p}$}] {$w_1$};
    \draw[reflexive above] (w1) to (w1);
    \node[world] (w2) [label={right:\mTrue{p}}, label={below:
        $\mSat{{}}{\Diamond\lnot p}$\\
        $\mSat/{{}}{\Box\Diamond\lnot p}$}, right=of w1] {$w_2$} ;
    \draw[reflexive above] (w2) to (w2);
    \node[world] (w3) [label={right:\mFalse{p}},right=of w2] {$w_3$};
    \draw[reflexive above] (w3) to (w3);
    \draw[->, bend left] (w1) to (w2);
    \draw[->, bend left] (w2) to (w3);
    \draw[->, bend left] (w3) to (w2);
    \draw[->, bend left] (w2) to (w1);
  \end{tikzpicture}
  \caption{\ref{nml:prf:dis:thm:KTBnot45} साठीचे प्रतिमान.}
  \label{nml:prf:dis:fig:KTBnot45}
\end{figure}

\begin{thm}\label{nml:prf:dis:thm:KD5not4}
  $\Log{KD5} \neq \Log{KT4} = \Log{S4}$.
\end{thm}

\begin{proof}
  \ref{nml:frd:acc:thm:soundschemas} वरून
  \Ax{D} आणि~\Ax{5} यांचे सर्व दृष्टांत
  प्रत्येक सर्वत्र उत्तरवर्ती युक्लिडी प्रतिमानात
  सत्य असतात. म्हणून असे प्रतिमान शोधणे पुरेसे
  आहे की त्यातील एका जगात~\Ax{4} चा एखादा
  दृष्टांत असत्य ठरेल. \ref{nml:prf:dis:fig:KD5not4}
  मधील प्रतिमान पाहा; येथे
  $\mSat/{M}{\Box p
  \lif \Box\Box p}[w_1]$.
\end{proof}

\begin{figure}[t]
  \centering
  \begin{tikzpicture}[modal]
    \node[world] (w2) [label=north west:\mTrue{p}]{$w_2$} ;
    \draw[reflexive left] (w2) to (w2);
    \node[world] (w1) [label=right:\mFalse{p},
      label=below:{$\mSat{{}}{\Box p}, \mSat/{{}}{\Box\Box p}$},
      below right=of w2]{$w_1$};
    \node[world] (w3) [label=north east:\mTrue{p}, above right=of w1]{$w_3$} ;
    \draw[reflexive right] (w3) to (w3);
    \node[world] (w4) [label=right:\mFalse{p},above right=of w2]{$w_4$};
    \draw[reflexive above] (w4) to (w4);
    \draw[->] (w1) to (w2);
    \draw[->] (w1) to (w3);
    \draw[->, bend left=15] (w2) to (w3);
    \draw[->, bend left=15] (w3) to (w2);
    \draw[->, bend left=15] (w2) to (w4);
    \draw[->, bend left=15] (w4) to (w2);
    \draw[->, bend left=15] (w3) to (w4);
    \draw[->, bend left=15] (w4) to (w3);
  \end{tikzpicture}
  \caption{\ref{nml:prf:dis:thm:KD5not4} साठीचे प्रतिमान.}
  \label{nml:prf:dis:fig:KD5not4}
\end{figure}

\begin{prob}
  $3$ जगांचे प्रतिमान वापरून
  \ref{nml:prf:dis:thm:KD5not4} ची पर्यायी
  सिद्धता द्या.
\end{prob}

\begin{prob}
  $\Log{KT4} \Proves/ \Ax{B}$ आणि
  $\Log{KT4} \Proves/ \Ax{5}$ हे दोन्ही दाखवणारे
  एकच परावर्ती संक्रमक प्रतिमान द्या.
\end{prob}













\section{संचातील सूत्र गृहीत धरून निष्पन्नता}\label{nml:prf:prg:sec}

\ref{nml:prf:prf:sec} मध्ये प्रणाली~$\Sigma$ मधील
सूत्राच्या सिद्धतेची संकल्पना परिभाषित केली.
आता संच~$\Gamma$ मधील सूत्रांपासून
$\Sigma$ मध्ये सिद्धतेसाठी ती विस्तारित करू.

\begin{defn}\label{nml:prf:prg:defn:Gammaproves}
  संच~$\Gamma$ मधील सूत्रांपासून
  प्रणाली~$\Sigma$ मध्ये सूत्र~$\varphi$
  हे निष्पन्न करता येण्याजोगे आहे, असे तेव्हा आणि तेव्हाच
  म्हणतो व $\Gamma \Proves[\Sigma] \varphi$ असे लिहितो,
  जेव्हा $n \geq 0$ आणि $\psi_1$, \dots, $\psi_n \in \Gamma$ अशी
  सूत्रे असतात की
  $\Sigma \Proves \psi_1 \lif (\psi_2 \lif \cdots (\psi_n \lif \varphi)
  \cdots)$.
येथे $n=0$ असताना ही implication मालिका फक्त $\varphi$ असते;
म्हणून कोणतीही गृहीतके न वापरता $\Sigma \Proves \varphi$ मिळते.
\end{defn}












\section{निष्पन्नतेचे गुणधर्म}\label{nml:prf:prp:sec}

\begin{prop}\label{nml:prf:prp:prop:derivabilityfacts}
  $\Sigma$ ही मोडल प्रणाली आणि $\Gamma$ हा मोडल
  सूत्रांचा संच असू द्या. पुढील गुणधर्म लागू होतात:
  \begin{enumerate}
  \item\label{nml:prf:prp:prop:derivabilityfacts-monotonicity} \emph{एकस्वनिकता}:
    $\Gamma \Proves[\Sigma] \varphi$ आणि
    $\Gamma \subseteq \Delta$ असेल, तर
    $\Delta \Proves[\Sigma] \varphi$;
  \item\label{nml:prf:prp:prop:derivabilityfacts-reflexivity}
    \emph{परावर्तिता}: $\varphi \in \Gamma$ असेल, तर $\Gamma
    \Proves[\Sigma] \varphi$;
  \item\label{nml:prf:prp:prop:derivabilityfacts-cut} \emph{कट}: $\Gamma
    \Proves[\Sigma] \varphi$ आणि $\Delta \cup \{\varphi\} \Proves[\Sigma] \psi$
    असेल, तर $\Gamma \cup \Delta \Proves[\Sigma] \psi$;
  \item\label{nml:prf:prp:prop:derivabilityfacts-deduction} \emph{निगमन प्रमेय}: $\Gamma \cup \{\psi\}
    \Proves[\Sigma] \varphi$ असेल तेव्हा आणि तेव्हाच $\Gamma \Proves[\Sigma] \psi \lif
      \varphi$;
  \item \label{nml:prf:prp:prop:derivabilityfacts-ruleT} \emph{नियम T}: जर
    $\Gamma \Proves[\Sigma] \varphi_1$ आणि \dots आणि $\Gamma
    \Proves[\Sigma] \varphi_n$ असेल आणि $\varphi_1 \to (\varphi_2 \lif \cdots (\varphi_n \lif
    \psi)\cdots)$ हा सर्वतः-सत्य दृष्टांत असेल, तर $\Gamma
    \Proves[\Sigma] \psi$.
  \end{enumerate}
\end{prop}

सिद्धता हा सोपा सराव आहे.
\ref{nml:prf:prp:prop:derivabilityfacts} मधील
\ref{nml:prf:prp:prop:derivabilityfacts-ruleT} या भागावरून,
उदाहरणार्थ, $\Gamma \Proves[\Sigma] \varphi \lor \psi$
आणि $\Gamma \Proves[\Sigma] \lnot \varphi$ असेल, तर
$\Gamma \Proves[\Sigma] \psi$ मिळते.
तसेच पुढे $\Gamma \cup \{ \varphi\} \Proves[\Sigma] \psi$
ऐवजी $\Gamma, \varphi \Proves[\Sigma] \psi$ असे लिहू.

\begin{defn}
  $\Gamma \Proves[\Sigma] \varphi$ असेल तेव्हा नेहमी
  $\varphi \in \Gamma$ असेल, तर आणि तेव्हाच संच
  $\Gamma$ हा प्रणाली~$\Sigma$ च्या सापेक्ष
  \emph{निगामीदृष्ट्या बंद} आहे, असे म्हणतो.
\end{defn}













\section{सुसंगतता}\label{nml:prf:con:sec}

सूत्रांच्या संचांचा सुसंगतता हा महत्त्वाचा
गुणधर्म आहे. सूत्रांच्या संचापासून
$\lfalse$ सारखा व्याघात निष्पन्न करता येण्याजोगे असेल,
तर तो संच विसंगत असतो; अन्यथा सुसंगत असतो.
संबंधित प्रणालीच्या सर्व स्वयंसिद्धकांचे सर्व आदेशन-दृष्टांत
प्रतिमानाच्या सर्व जगांत सत्य असतील, तर त्या प्रणालीच्या सापेक्ष
विसंगत संचातील सर्व सूत्रे प्रतिमानाच्या एका जगात
एकत्र सत्य असू शकत नाहीत. संपूर्णता प्रमेयासाठी आपण याचा
व्यत्यास सिद्ध करू: प्रत्येक सुसंगत संच
एका प्रतिमानातील एका जगात सत्य असतो; ते
प्रतिमान म्हणजे ``कॅनॉनिकल प्रतिमान.''

\begin{defn}
  संच~$\Gamma$ हा प्रणाली~$\Sigma$ च्या सापेक्ष
  \emph{सुसंगत}, किंवा आपण म्हणू त्याप्रमाणे
  $\Sigma$-सुसंगत, असेल तेव्हा आणि तेव्हाच
  $\Gamma \Proves/[\Sigma] \lfalse$.
\end{defn}

उदाहरणार्थ, $\{ \Box(p \lif q), \Box p, \lnot\Box q \}$
हा संच विधानीय तर्कशास्त्राच्या सापेक्ष सुसंगत आहे,
पण \Log{K}-सुसंगत नाही. तसेच
$\{ \Diamond p, \Box\Diamond
p \lif q, \lnot q \}$ हा संच \Log{K5}-सुसंगत नाही.

\begin{prop}\label{nml:prf:con:prop:consistencyfacts}
  $\Gamma$ हा सूत्रांचा संच असू द्या. मग:
  \begin{enumerate}
  \item $\Gamma$ हा $\Sigma$-सुसंगत असेल तेव्हा आणि
    तेव्हाच असे एखादे सूत्र~$\varphi$ असते की
    $\Gamma \Proves/[\Sigma] \varphi$.
  \item \label{nml:prf:con:prop:consistencyfacts-b}
    $\Gamma \Proves[\Sigma] \varphi$ असेल तेव्हा आणि
    तेव्हाच $\Gamma \cup \{
    \lnot\varphi \}$ हा $\Sigma$-सुसंगत नाही.
  \item \label{nml:prf:con:prop:consistencyfacts-c}
    $\Gamma$ हा $\Sigma$-सुसंगत असेल, तर
    कोणत्याही सूत्र~$\varphi$ साठी
    $\Gamma \cup \{ \varphi \}$ हा
    $\Sigma$-सुसंगत असतो किंवा
    $\Gamma \cup \{ \lnot\varphi \}$ हा
    $\Sigma$-सुसंगत असतो.
  \end{enumerate}
\end{prop}

\begin{proof}
  ही तथ्ये अभिजात विधानीय तर्कशास्त्रावरून सहज
  मिळतात. \ref{nml:prf:con:prop:consistencyfacts-c} साठी
  युक्तिवाद देऊ. प्रतिपरिवर्तनाने सिद्ध करू:
  $\Gamma \cup \{ \varphi \}$ आणि
  $\Gamma \cup \{ \lnot\varphi \}$ यांपैकी कोणताही
  संच $\Sigma$-सुसंगत नाही, असे समजा. मग
  सुसंगततेच्या व्याख्येवरून
  $\Gamma, \varphi \Proves[\Sigma]
  \lfalse$ आणि $\Gamma, \lnot \varphi \Proves[\Sigma] \lfalse$
  हे दोन्ही मिळतात. निगमन प्रमेयावरून
  $\Gamma \Proves[\Sigma] \varphi \to \lfalse$ आणि
  $\Gamma \Proves[\Sigma] \lnot\varphi \lif \lfalse$.
  पण $(\varphi \lif
  \lfalse) \lif ((\lnot\varphi \lif \lfalse) \lif \lfalse)$
  हा सर्वतः-सत्य दृष्टांत आहे; म्हणून
  \ref{nml:prf:prp:prop:derivabilityfacts}\ref{nml:prf:prp:prop:derivabilityfacts-ruleT}
  वरून $\Gamma \Proves[\Sigma] \lfalse$.
\end{proof}



